% ============================================================
\chapter{Conclusions and Future Work}
\label{chap:conclusion}
% ============================================================

\section{Summary of findings}
\label{sec:summary}

This thesis asked one question: when the correlation matrix in
hierarchical portfolio allocation is replaced by a causal graph, is
there a performance improvement, and how much of any improvement is
attributable to the graph's skeleton
versus its edge orientations? It made the question answerable by building
the
machinery it presupposes (a reproducible rolling discovery pipeline that
delivers the identical asset--asset DAG to every allocator, A1, and an
orientation-aware allocator family with skeleton and correlation controls
verified orientation-blind under stated invariances, A2, discharging
Requirement 2) and answered
it under multiplicity-, distribution- and stability-aware inference (A3).

The answer, in summary: the causal graph's gain over the
correlation matrix is stable across estimation windows, but its source
shifts: against the primary skeleton control, the skeleton pays most
clearly near the conventional one-year window, while orientation pays at
every estimation window and, robustly to the choice of orientation-blind control,
from one year out. The between-window differences are about one standard
error, and the window grid was widened partway through, so the shapes
carry that uncertainty. Orientation buys robustness to the estimation
window, not level, and a direction-free de-factored covariance
reproduces that robustness, so its mechanism is market-factor removal
rather than directional content. Family-wise, nothing rejects over the
full allocator grid of every strategy the experiment plan evaluated;
support reaches conventional significance
only under the narrower reported family, at the longest estimation windows,
marginally, where the point estimates also peak. In detail (net annualised Sharpe, 2007--2024,
monthly
rebalancing, 5\,bps costs, windows of 189/252/378/504 trading days):

\begin{enumerate}[itemsep=2pt]
  \item The total gain is stable across estimation windows.
        Replacing the correlation matrix with the DYNOTEARS graph is
        worth $+0.018$ to $+0.032$ net Sharpe (\skelsem{} $-$ \HRP{})
        at every window, and 36 of 40 graph-versus-control cells are
        positive.
  \item Its composition shifts with the window. The skeleton's
        contribution is hump-shaped and peaks at the one-year
        convention; orientation's is positive everywhere and U-shaped,
        largest where the sample moments are weakest. Orientation buys
        robustness to the estimation window, not level, and at the
        18-month window \orientsem{} attains the highest Deflated
        Sharpe ratio of all \rsNtrials\ configurations evaluated
        ($\rsDtwosDSRWthree$).
  \item The mechanism needs no graph; the skeleton does. A
        de-factored covariance with no graph reproduces the
        covariance-channel gain at every window, so the orientation
        robustness works through market-factor removal; the matching
        graph-free skeleton control fails to reproduce the skeleton's
        one-year gain, so the channel that pays at the conventional
        window does need the discovered graph
        (Section~\ref{sec:orientation-role}). The shapes themselves
        are anchor-dependent at nine months (limitation (vii)).
  \item Family-wise support depends on the family, and the
        full grid does not reject. Every SPA test runs over both the
        full seven-allocator grid and the narrowed five-allocator
        family fixed after the results were seen: the full grid rejects
        nowhere, the narrowed family rejects marginally at the longest
        estimation windows only, and the pooled test that prices the
        strategy and window searches jointly rejects for neither family
        (Table~\ref{tab:spa}). The $90\%$ Model Confidence Set declines
        to separate treatments from controls.
  \item The effect is method- and allocator-specific. Under
        VARLiNGAM graphs orientation adds nothing at any window
        (limitation (ii)); under the magnitude-degraded ridge-Granger
        control the ordering survives but the SEM covariance collapses,
        so the value sits in trustworthy edge magnitudes; and on a
        second family member (HERC) only orientation's sign generalises
        (Section~\ref{sec:decomposition}). Descriptively, the edge
        concentrates in calm regimes; the stress hypothesis is
        rejected.
\end{enumerate}

\section{Contributions}
\label{sec:contributions}

The contribution of this thesis is the orientation-aware allocator family
and the fixed-graph ablation design; the decomposition of
Chapter~\ref{chap:evaluation} is what that instrument found. That single
claim unpacks into a methodological core and two strands that flow from
it:

\begin{itemize}[itemsep=2pt]
  \item Methodological (the core): the orientation-aware allocator family
        (SEM-implied, de-standardised, residual-calibrated
        covariance allocation, and causal-ordered bisection, which replaces
        the clustering dendrogram with the DAG's topological order),
        together
        with the fixed-graph ablation design that makes skeleton and
        orientation separately measurable, with controls verified
        orientation-blind under stated invariances. To my knowledge, both the family and the decomposition are
        new; the closest prior work \cite{howard2025causal} never sets
        weights from its graphs.
  \item Empirical (what the instrument found): the decomposition itself, traced across four
        estimation windows (a stable total; skeleton hump-shaped and
        peaking at the one-year convention against the primary anchor;
        orientation positive at every
        estimation window against that anchor and robustly so from one year out;
        family-wise support absent under the full grid and
        marginal, confined to the longest estimation windows, under the narrowed
        family), the identification of its covariance-channel mechanism
        as market-factor removal by a direct control, and its
        method- and
        regime-dependence.
  \item Meta-scientific: a
        reproducible, fully deterministic pipeline (including a
        data-determinism defect
        whose fix overturned two previously-drawn conclusions), with
        every design extension disclosed as such, all cells
        cache-verified, and the
        harness gated on an exact replication of the prior phase's
        result.
\end{itemize}

\section{Limitations}
\label{sec:limitations}

(i) The decomposition's components are
$0.01$--$0.03$ Sharpe: near the 18-year detection floor, individually
non-significant against the stronger controls, and quantifiably
underpowered: the one-sided minimum detectable effects at $80\%$ power
are ${\approx}\,\rsContrastMDETotal$ (total) and
${\approx}\,\rsContrastMDEOrient$ (orientation) net Sharpe, roughly
twice the observed increments (Section~\ref{sec:believe}).
The point estimates are exactly reproducible; the inference, where it
clears the family-wise bar, is marginal, clears it only under the
post-hoc narrowed family, not under the full grid, and does not survive
pricing the window search jointly with the strategy search.
(ii) The
orientation effect exists under one discovery method (DYNOTEARS); until
another method reproduces it, it should be read as a property of DYNOTEARS'
contemporaneous structure, not of markets, and the de-factored control
shows its covariance channel is reproducible with no graph at all
(the skeleton channel's one-year gain, by the matching control, is
not). (iii) Long-only US
large-cap equity only; every variant absorbs the full ${\approx}50\%$ GFC
drawdown; transfer to richer universes is untested. (iv) The
universe approximates the S\&P~100 by market-cap ranking and a fixed
snapshot union, with mild residual survivorship effects. (v) The
asset-block SEM marginalises over the macro drivers, so the structural
shocks' assumed independence is an approximation. (vi) The choice
of research question was informed by the first phase's
results, the window grid was widened to four estimation windows partway through
the study, and the family-wise comparison set was narrowed from the
full seven-allocator grid to five after the results were seen
(Section~\ref{sec:believe}); all three are disclosed, both families'
numbers are reported, and the unified
\rsNtrials-trial deflation universe prices in every configuration the
search ever touched, but a deflation correction is not a substitute for
out-of-sample replication. The 2025--26 slice
(Section~\ref{sec:oos}), whose protocol was fixed before the data
existed, supplies the directional part of that
replication: every sign it specified came back
consistent; at 19 months
it can certify signs and nothing more, so the significance question
stays open. (vii) At the shortest window the two
orientation-blind controls themselves disagree (\skelalt{} $-$ \skelsamp{} $= +0.025$,
$p=0.013$), so the nine-month contributions carry extra estimator
uncertainty;
more generally the hump and U shapes of the decomposition are anchored on
\skelsamp{}, and under the \skelalt{} anchor the skeleton contribution is
flat
and the nine-month orientation contribution slightly negative
(Section~\ref{sec:decomposition}). (viii) Month-to-month
persistence of the discovered edges was not quantified: graph stability
over time is attested by the determinism of the fits and the stability of
aggregate diagnostics, not by an edge-overlap measure.

\section{Broader impact}
\label{sec:impact}

\paragraph{For investors.} Allocation research is consumed by asset
managers
and pension schemes, so its failures are borne, diffusely and years
later, by end savers. The characteristic harm in this literature is not a
malfunctioning system but a convincing backtest of an overfit strategy;
multiplicity-corrected, family-wise-honest evaluation is therefore not
methodological decoration but the investor-protection function of research
method, and this report applies it to its own headline. The same honesty is
a corrective to what might be called causal-washing: a ``causal''
label invites over-trust, and the finding here is that a causal method's
value can be modest, window- and method-dependent, and, at the
conventional one-year window, largely attributable to structure that
cheap symmetrisation already captures. The label guarantees nothing; the
decomposition, not the branding, is what carries information.

\paragraph{For the field.} One of this project's technical choices
generalises into a norm. The determinism episode of Section~\ref{sec:cache}
(two conclusions
overturned by a reproducibility fix) argues that reproducibility
belongs among publication prerequisites rather than aspirations; the
deterministic, seed-free construction of the allocator family is the
design corollary, and it is what makes every number in this report
exactly replicable.

\paragraph{Costs.} The content-keyed cache cut the marginal cost of a
configuration from ${\sim}15$ hours to ${\sim}40$ seconds, which is what
made an \rsNtrials-configuration study, and hence its multiplicity
correction, affordable at all. The compute and energy economics of
caching-first experimental design are a transferable, if unglamorous,
contribution to sustainable empirical research: rigour and efficiency were
complements here, not a trade-off.

\section{Future work}
\label{sec:future}

(1) A second orientation source. The decomposition's biggest caveat
is method-dependence: repeating the fixed-graph ablation on graphs from a
different family (e.g.\ the non-linear NTS-NOTEARS, whose prior-transfer was
verified in a reduced-scope probe but whose ${\sim}10\times$ fit cost
deferred a full backtest) would test whether orientation value is a
DYNOTEARS artefact. (2)
Signed structural allocation. Equalising risk contributions over
structural shock origins (where risk enters the causal system)
rather than over correlated returns is the natural next use of
$\Sigma_{\mathrm{SEM}}$; a long-only variant was among the seven
allocators of the pre-registered grid and is priced into the full-grid
tests of Section~\ref{sec:believe-rigour}, but exact per-shock parity
generally requires
shorting, so a long--short mandate, not the long-only setting studied here,
is its natural home. (3) Richer universes, where cross-asset
causal chains are plausibly longer-lived than among 99 co-moving
large-caps:
multi-asset-class portfolios are where the skeleton/orientation split
could materially differ. (4) Orientation-aware risk attribution.
The
total-effect matrix $B$ supports ex-post decompositions, attributing a
month's loss to the shocks that drove it, that the correlation toolkit
cannot express;
worthwhile even where the allocation gain is modest. The
infrastructure built here (the deterministic pipeline, the discovery cache
and the ablation harness) makes each of these a matter of days, not
months, of compute, which, as this project found twice, is a prerequisite
for trusting the answers.
